\scanpage{191}\mspage{47}
% archive sheet number: 191

Dém. : a) \inserted{F1)} est évident, car $0$-connexe $\Rightarrow$ non vide. \inserted{F0) car $W_0$-coasph. $\Rightarrow$ $W_0$-équiv. i.e. $\pi_0(u)$ bij.}

Pour \uline{prouver b)}, il suffit de prouver que F1, F2 (si $I'$ est fblt filtrante) implique que les $\,{}_{i\backslash}I'$ sont fblt filtrantes (a fortiori $0$-connexes) \inserted{et de plus \deleted{\unsure{le fait}} que $I$ \unsure{l'est} également}. On sait déjà qu'ils sont non vides par F1, reste à voir que si $(i',\alpha),(j',\beta)$ sont dans $\,{}_{i\backslash}I'$, alors ils sont ``majorés'' par un $3^{\mathrm{o}}$ objet de $\,{}_{i\backslash}I'$.

\diagram{\begin{tikzcd}[ampersand replacement=\&,row sep=1.3em,column sep=1.5em]
\& i'\arrow[rr,"r"]\arrow[dd,no head] \& \& k'\arrow[r,dashed,"t"]\arrow[dd] \& l'\arrow[dd,no head]\\
\& \& j'\arrow[ur,"s"]\arrow[dd,no head,crossing over] \& \&\\
\& u(i')\arrow[rr,"{u(r)}"] \& \& u(k')\arrow[r,dashed,"{u(t)}"] \& \deleted{\unsure{u(l')u(l')|u(l'),u(l')}}\\
i\arrow[ur,"\alpha"]\arrow[rr,"\beta"'] \& \& u(j')\arrow[ur,"{u(s)}"'] \& \&
\end{tikzcd}}

Comme $I'$ est fblt filtr., on trouve

\diagram{\begin{tikzcd}[ampersand replacement=\&,row sep=small]
i'\arrow[rr,"r"] \& \& k'\\
\& j'\arrow[ur,"s"'] \&
\end{tikzcd}$,$}

\deleted{\unsure{d'où}} \deleted{\unsure{les deux flèches}} \inserted{$u(r),u(s)$} de $I$ (cf diagramme), appliquons F2 : la double-flèche $u(r)\alpha,u(s)\beta:i\rightrightarrows u(k')$, on trouve $k'\xrightarrow{t}l'$ (en pointillés sur diagramme) telle que $(u(t)u(r))\alpha=(u(t)u(s))\beta$, \inserted{soit $\gamma$,} ce qui montre que l'objet $(l',\gamma:\unsure{i'|i}\longrightarrow u(l'))$ majore $(i',\alpha)$ et $(j',\beta)$.

\inserted{Si $I'$ est fblt filtrant et $u$ satisfait F1, alors $I$ est fblt filtrant, c'est immédiat.}

\uuline{Prouvons c).} Il \unsure{reste} à prouver que si $I'$ est \uline{filtrante} (non seulement fblt filtrante), \deleted{alors} \deleted{\unsure{al}} \inserted{si} $u$ cofinal, alors F2 est vérifiée, et \deleted{que $I$} \inserted{de plus} \deleted{et plus} $\,{}_{i\backslash}I'$ \inserted{et $I$} sont filtrantes (pas seulement fblt. filtrantes). \deleted{\unsure{Or}} Mais notons \deleted{que} tout de suite que $\,{}_{i\backslash}I'\longrightarrow I'$ est coétale, donc $I'$ filtrante \inserted{donc loc. filtr.-te} implique que $\,{}_{i\backslash}I'$ est aussi loc. filtrante ; comme par hypothèse les $\,{}_{i\backslash}I'$ sont $0$-connexes, elles sont donc filtr.-tes. La condition F2 \deleted{\unsure{soit}} \unsure{sort|suit} facilement à présent :
